\section{Objectif principal}

Concevoir, implémenter et déployer un pipeline distribué d'annotation automatique
d'images géospatiales par \gls{ia}, capable d'exploiter les ressources \gls{gpu}
du cluster \gls{kubernetes} de la HEIG-VD.

\section{Objectifs secondaires}

\begin{enumerate}
    \item \textbf{Évaluation du stockage objet} \\
          Analyser et évaluer les alternatives à \gls{minio} (RustFS, CEPH)
          en termes de performance, de licence et d'intégration avec l'infrastructure existante.

    \item \textbf{Traitement distribué des images} \\
          Mettre en place un pipeline \gls{ray} capable de distribuer le traitement
          des images sur plusieurs \glspl{gpu} en parallèle.

    \item \textbf{Optimisation du chargement des images} \\
          Intégrer \gls{zarr} pour la conversion et la lecture par tuiles des images,
          afin de réduire la consommation mémoire et d'améliorer les performances.

    \item \textbf{Persistance des résultats géospatiaux} \\
          Stocker les polygones issus de la segmentation dans une base de données
          \gls{postgresql} avec l'extension \gls{postgis}.

    \item \textbf{Intégration avec Label Studio} \\
          Exporter ou synchroniser les annotations produites automatiquement
          vers \gls{labelstudio} pour permettre la validation humaine.

    \item \textbf{Déploiement sur Kubernetes} \\
          Packager et déployer l'ensemble du pipeline sous forme de ressources
          \gls{kubernetes} (Pods, Jobs, ou Deployments).
\end{enumerate}

\section{Hors périmètre}

Les éléments suivants sont explicitement hors du périmètre de ce travail :

\begin{itemize}
    \item le réentraînement ou la modification du modèle SAM3 ;
    \item le développement d'une interface utilisateur de visualisation ;
    \item la gestion de la sécurité et des accès au cluster (supposée existante).
\end{itemize}
